\section*{Chapter \ref{ch:g11:titration} --- Titration}

\begin{solution}{exo:g11:titration:1}
The \omterm{def:g11:titration:titration}{titrant} is the solution of accurately known concentration, in the
burette; the \omterm{def:g11:titration:titration}{titrated solution} contains the species to be measured, in
the flask. At \omterm{def:g11:titration:equivalence}{equivalence} the \omterm{def:g11:titration:titration}{titrant} added and the titrated species are
in the proportions of the equation: both are used up.
\end{solution}

\begin{solution}{exo:g11:titration:2}
$n = 0.0200 \times 0.0125 = \qty{2.50e-4}{mol}$.
\end{solution}

\begin{solution}{exo:g11:titration:3}
$V_E = \qty{14.3}{mL}$; the iron(II) line starts at
\qty{14.3e-4}{mol}, that is \qty{1.43e-3}{mol}.
\end{solution}

\begin{solution}{exo:g11:titration:4}
The \omterm{def:g11:titration:titration}{titrant} is itself coloured (purple) and its product almost
colourless: the colour appears in the flask as soon as the permanganate
is in excess.
\end{solution}

\begin{solution}{exo:g11:titration:5}
5: the equation consumes 5 iron(II) \omterm{def:g9:ions:ion}{ions} per permanganate \omterm{def:g9:ions:ion}{ion}, and at
\omterm{def:g11:titration:equivalence}{equivalence} the \omterm{def:g7:chemical-reactions:reactant}{reactants} have been brought together in these
proportions.
\end{solution}

\begin{solution}{exo:g11:titration:6}
$n(\ce{MnO4^-}) = 0.0200 \times 0.00840 = \qty{1.68e-4}{mol}$;
$n(\ce{Fe^{2+}}) = 5 \times \num{1.68e-4} = \qty{8.40e-4}{mol}$;
$c = \num{8.40e-4} / 0.0100 = \qty{0.0840}{mol/L}$;
$C_m = 0.0840 \times 55.8 = \qty{4.69}{g/L}$.
\end{solution}

\begin{solution}{exo:g11:titration:7}
Ratio 1:1: $n = 0.0100 \times 0.0114 = \qty{1.14e-4}{mol}$ in
\qty{10.0}{mL}, so \qty{1.14e-3}{mol} in the \qty{100.0}{mL}, that is
$\num{1.14e-3} \times 176.0 = \qty{0.201}{g}$: about \qty{200}{mg}.
\end{solution}

\begin{solution}{exo:g11:titration:8}
Each drop would not react at once: the colour of the \omterm{def:g11:titration:titration}{titrant} would
linger before \omterm{def:g11:titration:equivalence}{equivalence}, and the end would be seen too early (or one
would have to wait minutes after each drop).
\end{solution}

\begin{solution}{exo:g11:titration:9}
Too large: she has gone past \omterm{def:g11:titration:equivalence}{equivalence}, adding permanganate in excess.
Near \omterm{def:g11:titration:equivalence}{equivalence} she should have added the \omterm{def:g11:titration:titration}{titrant} drop by drop,
swirling, and stopped at the first lasting faint pink.
\end{solution}

\begin{solution}{exo:g11:titration:10}
$n(\ce{Fe^{2+}}) = 5 \times 0.0200 \times 0.0130 = \qty{1.30e-3}{mol}$ in
\qty{10.0}{mL}: $c = \qty{0.130}{mol/L}$ for the diluted solution, and
$10 \times 0.130 = \qty{1.30}{mol/L}$ for the concentrated one.
\end{solution}

\begin{solution}{exo:g11:titration:11}
$0.0630 \times 55.8 = \qty{3.52}{g}$.
\end{solution}

\begin{solution}{exo:g11:titration:12}
$0.1 / 7.2 \approx \qty{1.4}{\percent}$, against
\qty{0.7}{\percent} for \qty{14.3}{mL}. Titrate a larger volume (or a
more concentrated sample) of the solution, or use a more dilute \omterm{def:g11:titration:titration}{titrant}.
\end{solution}

\begin{solution}{exo:g11:titration:13}
$n(\ce{MnO4^-}) = 0.0200 \times 0.0176 = \qty{3.52e-4}{mol}$;
$n(\ce{H2O2}) = \frac{5}{2} \times \num{3.52e-4} = \qty{8.80e-4}{mol}$ in
\qty{10.0}{mL}: \qty{0.0880}{mol/L} in the diluted solution,
\qty{0.880}{mol/L} in the disinfectant, that is
$0.880 \times 34.0 = \qty{29.9}{g/L}$ (about a 3 \% solution).
\end{solution}

\begin{solution}{exo:g11:titration:14}
$n(\ce{MnO4^-}) = \num{1.4e-3} / 5 = \qty{2.8e-4}{mol}$, to be delivered
in about \qty{0.015}{L}: $c \approx \qty{0.019}{mol/L}$, so a
\qty{0.0200}{mol/L} solution. Ten times more concentrated, $V_E$ would be
about \qty{1.4}{mL}: the burette error would reach some
\qty{7}{\percent}.
\end{solution}

\begin{solution}{exo:g11:titration:15}
$8 \times 0.0200 \times 0.0143 = \qty{2.29e-3}{mol}$ of hydrogen \omterm{def:g9:ions:ion}{ions}.
\qty{10}{mL} at \qty{1.0}{mol/L} hold \qty{1.0e-2}{mol}: enough, more
than four times over. Without acid the reaction of the chapter could not
take place as written (hydrogen \omterm{def:g9:ions:ion}{ions} are a \omterm{def:g7:chemical-reactions:reactant}{reactant}), and permanganate
would react differently.
\end{solution}

\begin{solution}{pb:g11:titration:1}
\textbf{1.} \ce{MnO4^-}/\ce{Mn^{2+}} and \ce{Fe^{3+}}/\ce{Fe^{2+}}. The
\omterm{def:g11:redox:oxidant}{oxidant} is the permanganate \omterm{def:g9:ions:ion}{ion}, the \omterm{def:g11:redox:oxidant}{reductant} the iron(II) \omterm{def:g9:ions:ion}{ion}.

\textbf{2.} \ce{MnO4^- + 8H+ + 5e- -> Mn^{2+} + 4H2O};
\ce{Fe^{2+} -> Fe^{3+} + e-}.

\textbf{3.} The second is multiplied by 5:
\ce{MnO4^- + 8H+ + 5Fe^{2+} -> Mn^{2+} + 5Fe^{3+} + 4H2O}. Charge on
the left $-1 + 8 + 10 = +17$; on the right $+2 + 15 = +17$.

\textbf{4.} Hydrogen \omterm{def:g9:ions:ion}{ions} are a \omterm{def:g7:chemical-reactions:reactant}{reactant}: the reaction needs an \omterm{def:g9:acids-bases-ph:acidic}{acidic
solution}.

\textbf{5.} It is total, fast and the only reaction of the permanganate
in the flask; the permanganate, purple, is used up before \omterm{def:g11:titration:equivalence}{equivalence} and
colours the flask pink after it.

\textbf{6.} In the burette; read at the bottom of the meniscus, eye
level, before and after, to \qty{0.05}{mL}.

\textbf{7.} Pale yellow (iron \omterm{def:g9:ions:ion}{ions}); each purple drop is used up at once
by the iron(II) \omterm{def:g9:ions:ion}{ions} still present.

\textbf{8.} $0.0200 \times 0.0143 = \qty{2.86e-4}{mol}$.

\textbf{9.} $n(\ce{Fe^{2+}}) = 5 \times \num{2.86e-4} =
\qty{1.43e-3}{mol}$.

\textbf{10.} $\num{1.43e-3} \times 55.8 = \qty{0.0798}{g} =
\qty{79.8}{mg}$.

\textbf{11.} $(80 - 79.8)/80 \approx \qty{0.3}{\percent}$.

\textbf{12.} $\qty{0.1}{mL}$ out of \qty{14.3}{mL} is \qty{0.7}{\percent},
about \qty{0.6}{mg}: the result, $79.8 \pm 0.6$ mg, agrees with the
label.

\textbf{13.} $14.1$: $5 \times 0.0200 \times 0.0141 \times 55.8 =
\qty{78.7}{mg}$; $14.4$: \qty{80.4}{mg}.

\textbf{14.} $(79.8 + 78.7 + 80.4)/3 = \qty{79.6}{mg}$.

\textbf{15.} The tablets themselves differ slightly in their iron
content, by about \qty{1}{\percent}.

\textbf{16.} No: $V_E$ would be about \qty{2.9}{mL}, and the burette
error \qty{3.5}{\percent}, five times worse.

\textbf{17.} Vitamin C would also be oxidised by the permanganate: the
reaction would no longer be unique, $V_E$ would be too large, and the
iron would be overestimated.

\textbf{18.} $\num{1.43e-3} \times 151.9 = \qty{0.217}{g}$, about
\qty{217}{mg}.

\textbf{19.} $\num{1.43e-3} \times \num{6.02e23} = \num{8.6e20}$ \omterm{def:g9:ions:ion}{ions}.

\textbf{20.} About \qty{80}{mg} of iron per tablet: the label is
honest.
\end{solution}
